% PINS ワークショップ論文（研究会原稿想定・6–8頁）
% ビルド: cd Research/paper && latexmk main.tex
% 体裁は暫定（jlreq 二段組）。投稿先決定後に ipsj.cls 等へ差し替える。
%
% 2026-07-31 改稿: 二者交渉era の論文からゲート型アーキテクチャの論文へ移植。
%   §4 = pins_gated_draft.md の 4.1--4.4（主張の本体）
%   §5 = 旧 主結果1--4 を圧縮（交渉機構の結果，副次）
\documentclass[uplatex,twocolumn,fontsize=9pt,paper=a4]{jlreq}
\usepackage[dvipdfmx]{graphicx}
\usepackage{booktabs}
\usepackage{amsmath}
\usepackage{url}

\newcommand{\todo}[1]{\textbf{【TODO: #1】}}
% 有意差マーク: 32シード対応比較で 95%CI が 0 を含まない
\newcommand{\sig}{$^{*}$}

\title{LLMはそのコストに見合うか：\\効率と保護を分離するゲート型GPU資源配分}
\author{Lay Kim Seng（東北大学 サイバーサイエンスセンター）\\ \todo{共著者・指導教員を確認}}
\date{}

\begin{document}
\maketitle

%-------------------------------------------------
\section*{概要}
高負荷のGPUクラスタにおける資源配分は「どのジョブの安全マージンを削るか」という
配給問題であり，近年はLLMに推論させる設計が提案されている．
本稿はその前提を問い直し，\textbf{LLMがどこでコストに見合い，どこで見合わないか}を
分解して測定する．Alibaba v2020 実トレース再生（32シード対応比較，2つの進捗法則）と
81件の事前登録済みテキスト例外スイートを用い，次を示す．
（1）数値だけの配給では\textbf{決定論的な入札・清算市場}が利用率・有用利用率・後悔の
全てで有意に勝ち，しかも\textbf{ゼロトークン}である．
（2）LLMの再現可能な寄与は\textbf{予約水準という単一スカラー}に収束し，
それ以外の毎ジョブ推論は純粋なコストである；保護と効率は同一資源の
\textbf{1つのダイヤル}であって，データが選ぶ勝者ではない．
（3）検証済みオークションを既定とし文脈トリガでのみ推論へ escalate する
\textbf{ゲート型アーキテクチャ}は効率中立であり，
テキストのない世界では\textbf{呼び出し0回・変更0件の厳密な no-op} となる．
（4）推論がコストに見合うのはスカラーで表現できない\textbf{テキスト例外}上であり，
予算を揃えた best-of-$N$ 単一呼び出し 29/81 に対し\textbf{討論}は 43/81
（片側正確McNemar $p=0.0007$）と有意に勝つ．
「必要になるまで何も要さない能力」という対が本稿の主結果である．

%-------------------------------------------------
\section{はじめに}
AI・科学技術計算ワークロードはバースト的かつ異質であり，静的なバッチスケジューラ
では資源利用率とサービス品質（SLA遵守率）の両立が難しい．利用率を押し上げると
SLA違反率も上がる——この緊張関係は，飽和領域では\textbf{全ジョブに安全マージンを
与えることが物理的に不可能}であることに由来する．すなわちサービス品質を実際に
決めているのは「\textbf{誰の安全マージンを削るか}」という配給の意思決定である．

この配給にLLMの推論を持ち込む設計は自然に見える．しかし\textbf{LLMは本当に
そのコストに見合っているのか}．本稿の立場は，この問いを全体としてではなく
\textbf{層ごとに分解して}測ることである．結論を先に述べると，分解の結果は
非対称であった：数値上の配給では決定論的機構が単独で，かつ無償で勝つ；
LLMの寄与は単一スカラーに縮退する；そして推論が明確に勝つのは，
数値が含まない\textbf{テキスト}として例外が到来する場面に限られる．

設計原則は一貫して「\textit{LLMは推論・説明し，決定論的機構が決定する}」であり，
本稿はその分業線を標語ではなく\textbf{測定された位置}として提示する．

主な貢献は以下の通りである．
\begin{enumerate}
  \item \textbf{効率は決定論的市場のもの}：明示的な入札・清算市場が床に対し
        利用率・有用利用率・後悔で有意に勝ち，従来最良であった交渉armを
        全6セル（プール$\times$法則）で上回る．コストはゼロトークンである（\S\ref{sec:market}）．
  \item \textbf{LLMの寄与は1スカラー，保護と効率は1ダイヤル}：予約水準のみを
        LLMに選ばせる \textbf{composed} armが完全なレフェリーarmを支配し，
        毎ジョブ推論は後悔と有用仕事を悪化させるだけである．保護は遊休GPUを
        持つことで買われ，それは市場の利用率利得そのものである（\S\ref{sec:scalar}）．
  \item \textbf{ゲート型アーキテクチャは効率中立}：検証器は正しい清算に対し
        ビット同一であり，escalation は SLA中立．テキストのない再生世界では
        トリガが843回発火して\textbf{呼び出し0・変更0}，割当はティック単位で
        市場と同一という\textbf{厳密な}境界が得られる（\S\ref{sec:gated}）．
  \item \textbf{推論が稼ぐ場所の特定}：事前登録81件のテキスト例外で，
        予算を揃えた best-of-$N$ に対し討論が 43/81 vs 29/81（$p=0.0007$）．
        予算だけでは何も買えない（29 vs 27, $p=0.250$）ことを対照が示す（\S\ref{sec:text}）．
  \item \textbf{特異性のコストの明示}：例外のない対照17件では市場 16/17 に対し
        全LLM arm が 12/17 である．既定で推論を起動することは routine 上で
        \textbf{有害}であり，これがトリガの実測的根拠である（\S\ref{sec:text}）．
\end{enumerate}

%-------------------------------------------------
\section{関連研究}
\paragraph{LLM駆動スケジューリング}
LLMsched \cite{llmsched} は単一LLMの提案をILPが修復する reason-then-guarantee 構成で
Kubernetes 比 JCT $-23.7\%$ 等を報告した．本研究はこの背骨を継承しつつ，
\textbf{推論層が何を買っているか}を分解して測る点で異なる．実測では，数値配給に
限れば推論層を除いた決定論的市場が効率で勝つ（\S\ref{sec:market}）．
\paragraph{学習ベーススケジューラ}
DeepRM/Decima 系のDRLは割当方策を学習するが，ブラックボックスであり分布シフトに
弱い．本研究は同一インタフェース上の表形式Q学習を弱い代表として測定し，
リターン分散が行動効果を支配する構造的困難を示す（\S\ref{sec:nego}）．
\paragraph{マレアビリティ}
実行時のジョブ伸縮は20年以上研究されつつ本番系では未活用である \cite{tarraf2024}．
本研究の弾力的部分割当はこの前提の上に立ち，その帰結（全か無かのEASYが不利になる）
も正直に報告する．
\paragraph{メカニズムデザイン}
均一価格オークションや予算制約による正直申告の誘導は古典的主題である．本研究の
貢献は，宣言的なLLMエージェントを含む閉ループでの実測（best-response試験）である．
\todo{各段落に文献を2--3本ずつ補充}

%-------------------------------------------------
\section{ゲート型アーキテクチャ}\label{sec:design}
\todo{図1: パイプライン図（入札・清算市場 → 検証器 →〔トリガ〕→ 討論エスカレーション）}

\subsection{全体像}
毎ティック，需要側の入札と供給側の売り気配を\textbf{決定論的に清算}して割当を得る．
清算結果は\textbf{検証器}が容量・重複・未知IDについて検査する．
検証を通り，かつ\textbf{文脈トリガ}が発火しなければ，その割当をそのまま適用する
（このパスはLLMを一切呼ばない）．トリガが発火したときに限り推論へ escalate する．

\begin{quote}\small
入札／売り気配 $\rightarrow$ 決定論的清算 $\rightarrow$ 自動検証\\
$\;$ 検証OKかつトリガ非発火 $\rightarrow$ オークションの割当を適用\\
$\;$ トリガ発火かつ裁定が有効 $\rightarrow$ エスカレーション arm の割当\\
$\;$ それ以外 $\rightarrow$ フォールバック（床）
\end{quote}

トリガは決定論的である：初回，直前の裁定が実行不能，空きGPUのバケット跨ぎ，
締切に対し新たに遅延したジョブ，prod層の新規到着．
\textbf{待ち時間がすでに発生している}ジョブにのみ推論コストを乗せる設計であり，
推論はクリティカルパスに乗らない（不収束時は床へフォールバックする）．

\subsection{エスカレーション arm：置換ではなく訂正}
エスカレーションの形状は本稿の結果を大きく左右する．市場の割当を破棄して
\textbf{全ジョブを再決定}する形状は prod層SLAを $-12.0$pt\sig 悪化させる
（\S\ref{sec:text}，Exp 84）．採用するのは\textbf{訂正型}である：市場が配分し，
レビュアは\textbf{符号つき差分}の提案のみを行い，しかも\textbf{注記を持つジョブ}
についてのみ提案できる．提案は決定論的に実行可能性検査され，通らなければ
市場の割当がそのまま残る（コードは決して裁定を修復しない）．

\subsection{二者交渉プロトコル（\S\ref{sec:nego}の対象）}
比較対象として，需要・供給エージェントが\textbf{有界譲歩プロトコル}で
マージンと予約を交換する交渉armを保持する．LLMの出力は「順序クラス・ヘッジ水準・
予約水準＋一行の正当化」に限定され，数値の最終決定は機構が行う．停止性は
ラウンド上限により保証される．
\todo{\texttt{method\_referee.tex}（162行）に毎ティック裁定型レフェリーの詳細が
あるが，\S\ref{sec:scalar} が referee を composed で置換する結論を出しているため，
そのまま採録せず「エスカレーション arm」節へ圧縮して統合すること}

\subsection{Stage-1: 使用量予測}
利用者申告は実使用をほぼ追跡しない（中央値50\% vs 0.22\%，$\rho=0.23$）．
PINSは申告をLLMに数値予測させる代わりに——LLMは較正された数値予測者として
一貫して敗北する——GBT分位点回帰で実使用量 $P10/P50/P90$ を予測し，
$P50$ を要求の基礎，区間幅を prod 層ヘッジ（newsvendor 的 prod-p90 規則）に用いる．

%-------------------------------------------------
\section{評価}\label{sec:eval}
\subsection{実験設定}
Alibaba cluster-trace-gpu-v2020 の606kジョブから連続窓を抽出し，到着・所要・需要を
実データのまま再生する（\texttt{trace\_replay}）．合成するのはトレースが持たない量
（層・締切・緊急度；シード固定のレシピ）のみである．プール$\{4,6,8\}$GPU，
1/4GPU量子，32シード対応比較（\sig は95\%CIが0を除く）．進捗は Amdahl 則と
飽和則の\textbf{両方}で実行し，\textbf{両法則で成立するときのみ}主張として数える．
LLMはローカルOllama（qwen2.5, gemma2, llama3）；比較arm間で窓とシードを完全対応させる．
基線：交渉なし規則（床），単一LLM，孤立エージェント，EASYバックフィル，表形式Q学習．

\subsection{決定論的市場が効率をゼロコストで取る}\label{sec:market}
推論層を明示的な入札・売り気配市場に置換する．床に対する \texttt{market}：

\begin{table}[t]
\centering\small
\caption{\texttt{market} $-$ 床（pool 8，32シード対応差，pt）}
\label{tab:market}
\begin{tabular}{lccccc}
\toprule
法則 & $\Delta$SLA & $\Delta$prod & $\Delta$util & $\Delta$有用 & $\Delta$後悔 \\
\midrule
amdahl & $-1.6$\sig & $-2.7$ & $+2.1$\sig & $+2.1$\sig & $-2.8$\sig \\
sat    & $-0.2$     & $-0.8$ & $+2.3$\sig & $+0.5$\sig & $-0.5$\sig \\
\bottomrule
\end{tabular}
\end{table}

従来の評価で勝っていた \texttt{negotiated} arm に対しても，市場は利用率・有用利用率・
後悔を\textbf{全6セル}（プール$\times$法則）で有意に改善する（例：amdahl pool 8 で
$\Delta$util $+1.1$\sig，$\Delta$有用 $+1.8$\sig，$\Delta$後悔 $-3.0$\sig）．
本プロジェクトの履歴上，有用利用率が正かつ後悔が負を同時に満たした唯一のarmである．
一方で prod層は保護しない（$\Delta$prodSLA は小さく全て非有意）．
そして\textbf{ゼロトークン・フォールバック0\%}である．効率は決定論的市場のものである．

\subsection{LLMの寄与は1スカラー，保護と効率は1ダイヤル}\label{sec:scalar}
レフェリーarmの唯一の再現可能な効果は prod層保護（$\Delta$prodSLA $-8.9$\sig/$-6.3$\sig）
である．そこで他の全てを切除する：\textbf{composed} arm はLLMに\textbf{予約水準のみ}を
選ばせ，残りのプールを市場に渡す（表\ref{tab:composed}）．

\begin{table*}[t]
\centering\small
\caption{床との差（pool 8，32シード対応差，pt）。composed は予約水準のみをLLMが選ぶarm}
\label{tab:composed}
\begin{tabular}{llccccc}
\toprule
法則 & arm & $\Delta$SLA & $\Delta$prod & $\Delta$util & $\Delta$有用 & $\Delta$後悔 \\
\midrule
amdahl & referee  & $-1.4$ & $\mathbf{-8.9}$\sig & $-1.2$ & $-1.9$\sig & $+3.2$\sig \\
amdahl & composed & $+0.4$ & $\mathbf{-5.6}$\sig & $-0.1$ & $-0.2$ & $+0.8$ \\
amdahl & market   & $-2.1$\sig & $-3.8$\sig & $+1.3$\sig & $+0.7$\sig & $-1.3$\sig \\
sat    & referee  & $-0.2$ & $\mathbf{-6.3}$\sig & $-0.8$ & $-2.2$\sig & $+4.4$\sig \\
sat    & composed & $-1.0$ & $\mathbf{-7.5}$\sig & $+0.2$ & $-0.8$ & $+1.7$\sig \\
sat    & market   & $+0.4$ & $-0.7$ & $+2.5$\sig & $+1.2$\sig & $-0.8$\sig \\
\bottomrule
\end{tabular}
\end{table*}

\paragraph{保護は切除に耐え，毎ジョブ推論は純コストだった．}
composed $-$ referee は prodSLA $+3.3$/$-1.1$（\textbf{両法則とも非有意}＝区別できない）
でありながら，後悔 $-2.4$\sig/$-2.6$\sig，有用利用率 $+1.7$\sig/$+1.4$ と改善する．
composed は referee を\textbf{支配し置換する}．完全なレフェリーの毎ジョブ陳述・反論・
裁定は，単一スカラーに対して何も買わず，後悔・有用仕事・24{,}240トークン/シードを費やす．

\paragraph{保護と効率は同一資源である．}
composed $-$ market では，予約が prod層保護を買う（飽和則で prodSLA $-6.7$\sig）一方，
util $-2.3$\sig，有用 $-2.0$\sig，後悔 $+2.5$\sig を\textbf{失う}．
予約はGPUを遊休させることで保護し，その遊休こそ市場の利用率利得の正体である．
これはデータが選ぶ勝者ではなく，運用者が選ぶ\textbf{1つのダイヤル}
（スループット対 prod層保護）である．

\paragraph{ワークロード上の交絡と，その主張への代償（Exp 96）．}
上記の生成器では層と締切スラックが単一の一様乱数から引かれており
（prod $\Leftrightarrow$ 緊急度 $\ge 1.667$ $\Leftrightarrow$ slack $\in[1.15,1.42]$，
点双列相関 $r=-0.78$），「層の保護」と「締切の厳しいジョブの保護」が分離できていなかった．
層ラベルを同一周辺分布のまま独立な乱数列から引き直し（$r=-0.02$，ジョブはビット同一），
両法則 $n=32$ で再実行した．結果は分裂する：amdahl では保護は減衰せず
（composed $-5.6$\sig $\rightarrow -3.4$\sig，差分の差 $+2.2\pm4.3$，非有意），
飽和則では多重比較を生き延びた唯一のセルが交絡解除に耐えない
（$-8.0$\sig, Holm $0.017$ $\rightarrow$ $-2.7$ 非有意, Holm $0.584$；
差分の差 $+5.3\pm5.8$, $p=0.053$）．正しい検定である差分の差は
どちらの法則でも有意でないため，保護の主張を撤回もせず，
きれいな層保護として言い換えもしない：\textbf{両方の世界の数値を併記し，
決着には約175対応シードを要すると述べる}．

同じ実験は明確な否定的結果も与える．厳しさ層（最厳三分位の違反率 \texttt{tight\_sla}）を
定義して層と脱相関させると，\textbf{どの推論armも最も締切の厳しいジョブを
どちらの法則でも保護しない}．飽和則ではレフェリーがそれらを有意に悪化させる
（$+6.9\pm5.1$\sig）一方，市場の当該層への効果は両世界でシード単位に同一である
——ラベルを一度も経由していないからである．機構内に厳しさを鍵にする要素はなく，
付与は凍結入札値の順である．これは\S\ref{sec:text}の特異性コストと同じ形であり，
最小スラック順という\textbf{未実行の決定論的レバー}を名指しする．

\subsection{ゲート型アーキテクチャは効率中立}\label{sec:gated}
\paragraph{検証器は正しい清算に対し不活性である．}
市場armを毎ティック検証器つきで再実行すると，両法則$\times$プール$\{4,6,8\}$$\times$32
シードで\S\ref{sec:market}の数値を\textbf{ビット同一に}再現し，フォールバックは全域0\%
である（Exp 86）．障害注入（過剰割当・負の付与・未知ID）では正しく棄却する．
整形式のオークション出力には決して触れない安全層——安全層はまさにそうあるべきである．

\paragraph{討論エスカレーションはオークション上でSLA中立である．}
文脈トリガの背後に討論＋運用マニュアルのエスカレーションを置く
（qwen2.5:14b, caps=predicted, pool 8, $n=32$）：

\begin{table}[t]
\centering\small
\caption{ゲート型アーキテクチャ（pool 8, $n=32$）。market は検証済みオークション単独，
gated はそれに討論を，corrected は\S\ref{sec:text}の訂正型エスカレーションを重ねたもの}
\label{tab:gated}
\begin{tabular}{lccc}
\toprule
arm & SLA & prodSLA & util \\
\midrule
床（LLMなし） & 53.9\% & 59.1\% & 80\% \\
market & 51.8\% & 55.3\% & 81\% \\
gated & 51.6\% & 54.7\% & 81\% \\
\textbf{corrected} & \textbf{51.8\%} & \textbf{55.3\%} & \textbf{81\%} \\
\bottomrule
\end{tabular}
\end{table}

gated $-$ market（対応差）は $\Delta$SLA $-0.20\pm1.40$（非有意），
$\Delta$prodSLA $-0.63$（非有意），$\Delta$util $-0.21$（非有意）．
トリガは全ティックの\textbf{15.7\%}で発火し（1076回，うち実行不能9），
52{,}853トークン/シードを消費した——エスカレーションは\textbf{実際に走り，
かつ結果を動かさなかった}．検証済みオークションの上に推論を重ねることは，
効率を損なわずに能力を追加する．

\texttt{corrected} 行は\S\ref{sec:text}でテキスト上勝つ\textbf{まさにその}
エスカレーションを載せた同一アーキテクチャである（Exp 92）．
32シードでトリガが843回発火しながら，素の市場と全指標で同一である
——統計的な引き分けではなく，ティック単位で同じ割当である．
これが効率中立性の主張を「非有意」ではなく\textbf{厳密}にしている．

\subsection{推論がコストに見合う場所：テキスト例外}\label{sec:text}
スカラーの予約は，内容が\textbf{テキスト}である例外——運用者の指示，
申告が誤りだと述べる注記，数値に現れない再起動——を表現できない．
事前登録した81件のテキスト依存例外で検証する．全armは同一の構造化
決定パケットを読む．比較は\textbf{構成上予算が揃っている}：単一呼び出しの横に，
討論と\textbf{同じ}呼び出し予算を与えた best-of-$N$ 単一呼び出しを走らせ，
2つの高コストarmの差を「予算の量」ではなく「予算の使い方」だけにする．

\begin{table}[t]
\centering\small
\caption{テキスト例外スイート（$n=81$，事前登録）}
\label{tab:text}
\begin{tabular}{lcc}
\toprule
arm & 処理成功 & LLM呼び出し \\
\midrule
market（数値のみ） & 0/81 & 0 \\
単一LLM & 27/81 & 81 \\
best-of-$N$（予算一致） & 29/81 & 569 \\
\textbf{討論} & \textbf{43/81} & 569 \\
\bottomrule
\end{tabular}
\end{table}

事前登録した主対比（討論 vs 予算一致 best-of-$N$）は
\textbf{$b=16$, $c=2$, 片側正確McNemar $p=0.0007$}（Exp 89）．
2つの対照が反証を難しくする．
\begin{itemize}
  \item \textbf{予算だけでは何も買えない．} best-of-$N$ は単一armの7倍の呼び出しを
        費やして2件しか増えない（29 vs 27, $p=0.250$, 非有意）．
        討論は\textbf{同一予算}で16件増やす．効果は支出ではなく\textbf{構造}に帰属する．
  \item \textbf{盲検の確認バッチが再現する．} 81件のうち50件は，先行する検出力不足の
        実行の\textbf{個別結果を見ずに}後から作成された．この50件のみで
        討論 29/50 vs best-of-$N$ 19/50（$b=11$, $c=1$, $p=0.0032$）．
\end{itemize}

機構はカテゴリ別の探索的分割に見える（事後・小セル・方向のみ読む）：討論の利得は
世界のモデルが\textbf{そもそも表現していない}例外に集中し（\textit{unmodeled}
13/21 vs 4/21，$b=9$, $c=0$；\textit{corrupt} 8/13 vs 5/13），
単一呼び出しが既に読める明示された方針では平坦である（\textit{nl\_policy} は
全LLM armで 8/21）．討論が効くのは\textbf{事実の欠落に気づく}必要があるときであり，
提示された指示に従うときではない．

\paragraph{推論すべきものが無い場所で推論するコスト．}
例外を含まない対照17件では，市場が 16/17 を得るのに対し\textbf{全LLM armが 12/17}
にとどまる．既定で起動される推論は routine 上で\textbf{積極的に有害}である．
これが\S\ref{sec:design}のトリガの実測的根拠である：エスカレーションは
テキストがある場所でのみコストに見合い，無い場所では走ってはならない．

\paragraph{正直な境界：同じ機構をテキストの無い場所で走らせると何も起きない．}
\S\ref{sec:text}への自然な反論は，作成されたスイートは実ワークロードについて
ほとんど何も示さない，というものである．そこで\textbf{この厳密に同じエスカレーション}を
——無改造・\texttt{--llm} 有効のまま——v2020再生のシミュレータ内で
\S\ref{sec:design}のトリガに繋いで実行した（Exp 92, pool 8, $n=32$）．
結果は表\ref{tab:gated}の \texttt{corrected} 行である：
市場と\textbf{全指標・全信頼区間・生き残った全Holm $p$値で同一}．
トリガは\textbf{843回}発火し，エスカレーションは毎回走り，
そして\textbf{LLM呼び出し0回・変更0件}であった：レビュア側も供給側もテキストで
ゲートされており，再生世界にはテキストチャネルが存在しない
（\texttt{Job} に注記フィールドは無く，トレースも記録しない——運用者の自由記述は
個人情報であり構成上除去されている）．

これは推論ではなく\textbf{厳密に述べられた境界}である：81件中43件のテキスト例外で
単一モデルの裁定を訂正する機構は，テキストが存在しない場所では\textbf{文字通りの
no-op} である．助けることができず，そして——配備上重要なのはこちらだが——
\textbf{害することも，コストを要することもできない}．

\paragraph{テキストの結果ではなくアーキテクチャの切除（Exp 84）．}
これとは\textbf{別の}，より侵襲的なエスカレーション——市場の割当を訂正するのではなく
テキストゲート無しで全ジョブを一から再決定するもの——は同じ世界で SLA中立かつ
prod層に有害である（$\Delta$SLA $+0.8$ 非有意，$\Delta$prodSLA $-12.0$\sig）．
これはそのまま報告する：\textbf{オークションの置換が有害である}という
\S\ref{sec:design}のゲート型アーキテクチャへの論拠であって，テキストについての
証拠ではない．2つの機構を混同すれば，アーキテクチャ上のコストを
チャネルの欠如に誤って帰属させることになった．

\paragraph{シミュレーション内SLA利得がなお手の届かない理由．}
公開されているGPUトレースはこの点をどちらにも決着させられない．
\S\ref{sec:gated}で escalate した15.7\%のティックは数値的な容量跨ぎであって
テキスト例外ではなく，市場はそれらを既に正しく裁いている——
つまりこのワークロードには，テキストを読む能力が改善すべきものが存在しない．
したがって\textbf{シミュレーション内のサービス品質利得は主張せず}，
かつ注記を持たないトレースに運用者注記を合成して作り出すことも\textbf{拒否する}．

本稿が行う主張は正確で完全に測定されている：ゲート型推論はテキストの無い場所で
\textbf{厳密に}効率中立であり（\S\ref{sec:gated}，Exp 92：呼び出し0・変更0・同一割当），
テキストのある場所で\textbf{正しい}（\S\ref{sec:text}，Exp 89：予算一致で 43/81 vs 29/81,
$p=0.0007$）．この対——\textbf{必要になるまで何も要さない能力}——は
捏造されたSLA利得より強く，かつ弁護可能であり，
ゲート型アーキテクチャが設計上果たすべきものそのものである．

%-------------------------------------------------
\section{交渉機構の結果（副次）}\label{sec:nego}
本節は\S\ref{sec:eval}に先立つ交渉era の測定であり，設計上の来歴として保持する．
主張の本体は\S\ref{sec:eval}である．

\paragraph{予測にもとづく適正要求と交渉の補完性．}
申告要求を Stage-1 $P50$ 使用量予測に置換すると，交渉@3bでSLA違反
$-9.2$〜$-17.0$pt\sig，飽和時 slowdown は半減する．予測要求の世界では
交渉@3bが床に対し\textbf{全ヘッドライン指標で同時に有意勝ち}する
（pool 4/6/8 で $\Delta$SLA $-4.9$/$-6.8$/$-6.2$\sig，
$\Delta$prodSLA $-5.7$/$-6.8$/$-5.2$\sig）．予測と交渉は代替ではなく補完である．

\paragraph{基線三角形．}
競合下（pool 6/8）で $\text{交渉@3b} \ge \text{床} > \{\text{EASY}, \text{Q学習}\}$．
EASYは予測実行時間で床に $+11.7$pt\sig 劣り，オラクル実行時間でも交渉に勝てない．
実行時間予測誤差はEASYに5--10ptのコストを課す——本システムで実行時間予測の質が
効く唯一のスロットである．表形式Q学習は同一インタフェース・同一情報でも床に
$+7.0$〜$+17.2$pt\sig 劣る：エピソードリターンの分散が行動効果を支配する．

\paragraph{機構がモデル規模を代替する．}
prod層保護は rule $\approx$ 3b $\approx$ 7b $\approx$ 14b で統計的に区別できず
（TOST等価），無ブレーキの単一LLMのみが規模を必要とし全規模で負ける．
能力の閾値は 0.5b と 1.5b の間にあり，gemma2:2b は族を跨いで通過，
llama3:8b は競合時に有意に悪化する——効くのは規模でなく命令追従性である．

\paragraph{インセンティブ層．}
無課金機構では虚偽申告が $-14$〜$-24$\sig の純利得を生む．利用者別予算＋係争ティック
課金の下で利得は全プールで統計的ゼロとなる（虚偽のコストは同一利用者のprodジョブに
$+15$\sig として着地する）．「自ジョブを最大化せよ」とだけ指示されたLLM利用者は
無課金機構を\textit{促されずに}搾取し，同じ課金で正直申告と等価になる．
抑止は\textbf{可読な規則}を通じて働く．

\paragraph{境界．}
（i）whole-GPU 割当は世界そのものに約10 SLA pt のパッキング損を課し，どの方策も
内部からは回収できない．（ii）第二トレース（MIT Supercloud）ではヘッジがプールの
25--100\%を占める粗さで交渉が害になる．（iii）交渉の全レバーはスラックを要する．

%-------------------------------------------------
\section{考察と限界}
\paragraph{分解が買うもの．}
LLMの寄与を単体の勝利としてではなく，Paretoダイヤル上の\textbf{1つの解釈可能な
スカラー}として報告することは，より正直であると同時により有用である：運用者は
ダイヤルを直接読み，供給側の目的をプロンプト1つの編集で差し替えられる（再学習不要）．
これがブラックボックスDRL方策に対する解釈可能性の優位を，測定された性質として
述べたものである．

\paragraph{SLAのレバーはマージンではなく受入である．}
本シミュレーションの競合下では，全体SLAを支配するのは\textit{どのジョブが始まるか}
であって\textit{どれが速く走るか}ではない．マージン1GPUは，プールが大半のジョブに
配れない加速である．したがって全体SLA上の推論利得には，マージンではなく
\textbf{受入や優先度}を動かすテキスト例外が要る．
\todo{Exp 63（受入レバー）：$n=32$ で実行中．H3（レフェリー vs 貪欲規則）は
32シード中2勝で\textbf{有意に負け}．H1/H2/H4は基線再実行の完了後に確定する}

\paragraph{妥当性への脅威．}
層・締切・緊急度は合成レシピであり，トレース再生世界は既存割当をプリエンプトしない
（自由なマレアビリティを仮定した主張はしない）．匿名課金は「全員が嘘つき」の世界を
救えない．LLM推論の実時間コストはゲートにより待ち時間の影に隠れる設計だが，
本稿の評価はシミュレーション時計上で行った．テキスト例外スイートは\textbf{作成された}
ものであり，運用からサンプルされたものではない——ただしこれは公開GPUトレースに
人手による例外テキストが存在しないことの帰結であり，合成による代替は
Exp 94 で不活性であることを確認した上で拒否している．
\todo{シミュレータ内のレフェリーは\S\ref{sec:text}の構造化パケットを使っていない
（\texttt{referee.py} は \texttt{packet.py} を参照しない）．
\S\ref{sec:nego}および受入レバーの否定的結果は，この\textbf{パケット以前の
インタフェース}上の結果として限定して述べること}

%-------------------------------------------------
\section{おわりに}
「LLMは推論・説明し，決定論的機構が決定する」という分業は標語ではなく，
全層で実測上最適な分割であった——そして本稿の分解はその線をより鋭くする．
数値上のGPU配給では，決定論的市場が単独で，かつゼロコストで\textbf{効率的に}決定し，
推論層の再現可能な寄与の全ては1次元Paretoダイヤル上の\textbf{単一の保護スカラー}に
縮退する．ゲート型アーキテクチャはその効率を無償で取り込み，文脈トリガ上でのみ
推論を起動する——それは効率中立である．推論層がコストに見合うのは，
スカラーが表現できない例外，すなわち\textbf{数値が含まないテキスト}の上であり，
そこではレビュアの討論が単一モデルの不確かな裁定を訂正する．
残された最前線は，そのテキストを実際に運ぶワークロードである：
運用者注記つきの実トレースの上で，\S\ref{sec:gated}の効率中立な能力が
測定可能なサービス品質の利得へ転化するはずである．

%-------------------------------------------------
\begin{thebibliography}{9}
\bibitem{llmsched} LLM-Driven Adaptive Cloud Resource Scheduling: Bridging Reasoning
Intelligence with Optimization Guarantees. IEEE Open J. Comput. Soc., 2026. \todo{正確な書誌}
\bibitem{tarraf2024} A.~Tarraf et al.: Malleability in Modern HPC Systems. IEEE TPDS, 2024. \todo{正確な書誌}
\bibitem{alibaba2020} Alibaba cluster-trace-gpu-v2020. \todo{書誌}
\bibitem{supercloud} MIT Supercloud Dataset. \todo{書誌}
\bibitem{easy} D.~Lifka: The ANL/IBM SP scheduling system (EASY). JSSPP, 1995. \todo{書誌}
\end{thebibliography}

\end{document}
